\documentclass[12pt, a4paper]{report}

% ─── Packages ───────────────────────────────────────────────────────────────
\usepackage[utf8]{inputenc}
\usepackage[T1]{fontenc}
\usepackage{babel}
\babelprovide[import, main]{indonesian}
\usepackage{geometry}
\usepackage{graphicx}
\usepackage{booktabs}
\usepackage{longtable}
\usepackage{multirow}
\usepackage{amsmath, amssymb}
\usepackage{hyperref}
\usepackage{caption}
\usepackage{subcaption}
\usepackage{float}
\usepackage{listings}
\usepackage{xcolor}
\usepackage{setspace}
\usepackage{indentfirst}
\usepackage{parskip}
\usepackage{titlesec}
\usepackage{fancyhdr}
\usepackage{cite}
\usepackage{url}
\usepackage{array}
\usepackage{tabularx}
\usepackage{pgfplots}
\pgfplotsset{compat=1.18}

% ─── Geometry ────────────────────────────────────────────────────────────────
\geometry{top=3cm, bottom=3cm, left=4cm, right=3cm}

% ─── Spacing ─────────────────────────────────────────────────────────────────
\onehalfspacing

% ─── Code Listing Style ──────────────────────────────────────────────────────
\lstdefinestyle{codestyle}{
  backgroundcolor=\color{gray!10},
  basicstyle=\ttfamily\footnotesize,
  breaklines=true,
  captionpos=b,
  commentstyle=\color{green!50!black},
  keywordstyle=\color{blue},
  stringstyle=\color{orange},
  numberstyle=\tiny\color{gray},
  numbers=left,
  frame=single,
  rulecolor=\color{gray!40},
  showstringspaces=false
}
\lstset{style=codestyle}

% ─── Header/Footer ───────────────────────────────────────────────────────────
\pagestyle{fancy}
\fancyhf{}
\fancyhead[R]{\leftmark}
\fancyfoot[C]{\thepage}
\renewcommand{\headrulewidth}{0.4pt}

% ─── Hyperref Setup ──────────────────────────────────────────────────────────
\hypersetup{
  colorlinks=true,
  linkcolor=black,
  citecolor=black,
  urlcolor=blue,
  pdftitle={Benchmarking PTQ dan QAT pada Arsitektur YOLO Modern untuk Inferensi Edge Heterogen},
  pdfauthor={Fatih Jawwad Al-Mumtaz}
}

% ════════════════════════════════════════════════════════════════════════════
\begin{document}

% ─── Cover ───────────────────────────────────────────────────────────────────
\begin{titlepage}
  \centering
  \vspace*{1cm}
  {\Large \textbf{BENCHMARKING POST-TRAINING QUANTIZATION DAN
  QUANTIZATION-AWARE TRAINING PADA ARSITEKTUR YOLO MODERN
  UNTUK INFERENSI EDGE HETEROGEN}}
  \vspace{1.5cm}

  {\large \textbf{SKRIPSI}}
  \vspace{1cm}

  Diajukan untuk memenuhi salah satu syarat memperoleh gelar\\
  Sarjana Teknik pada Program Studi Teknik Komputer\\
  Universitas Darussalam Gontor
  \vspace{1.5cm}

  % \includegraphics[width=4cm]{logo_unida.png}
  \vspace{1.5cm}

  Oleh:\\
  \textbf{Fatih Jawwad Al-Mumtaz}\\
  NIM: XXXXXXXXXX

  \vfill
  {\large
    PROGRAM STUDI TEKNIK KOMPUTER\\
    FAKULTAS SAINS DAN TEKNOLOGI\\
    UNIVERSITAS DARUSSALAM GONTOR\\
    PONOROGO\\
    \the\year
  }
\end{titlepage}

% ─── Front Matter ────────────────────────────────────────────────────────────
\pagenumbering{roman}
\setcounter{page}{2}

\chapter*{HALAMAN PERSETUJUAN}
\addcontentsline{toc}{chapter}{Halaman Persetujuan}
\vspace{2cm}
\noindent Skripsi berjudul \textit{``Benchmarking Post-Training Quantization dan
Quantization-Aware Training pada Arsitektur YOLO Modern untuk Inferensi
Edge Heterogen''} telah disetujui oleh pembimbing untuk diujikan.
\vspace{3cm}
\begin{flushright}
  Ponorogo, \ldots\ldots\ldots\ldots 2025\\[2cm]
  \underline{\hspace{5cm}}\\
  Dosen Pembimbing\\
  NIP: \ldots\ldots\ldots\ldots
\end{flushright}
\newpage

\chapter*{ABSTRAK}
\addcontentsline{toc}{chapter}{Abstrak}
\noindent\textbf{Fatih Jawwad Al-Mumtaz.}
\textit{Benchmarking Post-Training Quantization dan Quantization-Aware
Training pada Arsitektur YOLO Modern untuk Inferensi Edge Heterogen.}
Skripsi. Program Studi Teknik Komputer, Universitas Darussalam Gontor,
Ponorogo. \the\year.

\vspace{0.5cm}
Deployment model deteksi objek berbasis YOLO pada perangkat edge
dengan sumber daya terbatas memerlukan teknik optimasi inferensi yang
tepat. Dua pendekatan utama yang umum digunakan adalah
\textit{Post-Training Quantization} (PTQ) dan
\textit{Quantization-Aware Training} (QAT), yang masing-masing memiliki
trade-off antara akurasi, latensi, dan biaya komputasi.
Penelitian ini melakukan benchmarking sistematis terhadap PTQ dan QAT
pada arsitektur YOLOv8n dan YOLOv26n di tiga platform hardware heterogen:
Apple Silicon M4 (Mac Mini M4), GPU NVIDIA RTX 4060, dan MCU-class
ESP32-S3. Metrik yang diukur meliputi latensi (mean, p95, p99),
throughput (FPS), akurasi (mAP@0.5), ukuran model, dan konsumsi daya.
Hasil penelitian diharapkan dapat memberikan panduan praktis bagi
pengembang sistem edge AI dalam memilih teknik optimasi yang sesuai
dengan karakteristik hardware target.

\vspace{0.5cm}
\noindent\textbf{Kata kunci:} kuantisasi, YOLO, edge AI, inferensi heterogen, PTQ, QAT, Apple Silicon, ESP32
\newpage

\chapter*{ABSTRACT}
\addcontentsline{toc}{chapter}{Abstract}
\noindent\textbf{Fatih Jawwad Al-Mumtaz.}
\textit{Benchmarking Post-Training Quantization and Quantization-Aware
Training on Modern YOLO Architectures for Heterogeneous Edge Inference.}
Undergraduate Thesis. Department of Computer Engineering,
Universitas Darussalam Gontor, Ponorogo. \the\year.

\vspace{0.5cm}
Deploying YOLO-based object detection models on resource-constrained
edge devices requires careful selection of inference optimization
techniques. Post-Training Quantization (PTQ) and Quantization-Aware
Training (QAT) are the two dominant approaches, each presenting
distinct trade-offs between accuracy, latency, and computational cost.
This study conducts a systematic benchmark of PTQ and QAT on YOLOv8n
and YOLOv26n across three heterogeneous hardware platforms: Apple
Silicon M4 (Mac Mini M4), NVIDIA RTX 4060 GPU, and MCU-class ESP32-S3.
Evaluation metrics include latency (mean, p95, p99), throughput (FPS),
accuracy (mAP@0.5), model size, and power consumption. Findings are
expected to provide actionable guidance for edge AI developers in
selecting optimization strategies suited to their target hardware.

\vspace{0.5cm}
\noindent\textbf{Keywords:} quantization, YOLO, edge AI, heterogeneous inference, PTQ, QAT, Apple Silicon, ESP32
\newpage

\tableofcontents
\newpage
\listoffigures
\newpage
\listoftables
\newpage

% ─── Main Matter ─────────────────────────────────────────────────────────────
\pagenumbering{arabic}
\setcounter{page}{1}

% ════════════════════════════════════════════════════════════════════════════
\chapter{PENDAHULUAN}
% ════════════════════════════════════════════════════════════════════════════

\section{Latar Belakang}

Perkembangan model deteksi objek berbasis deep learning, khususnya
keluarga YOLO (\textit{You Only Look Once}), telah memungkinkan
penerapan computer vision secara real-time pada berbagai aplikasi
praktis seperti pengawasan, robotika, dan sistem otonom \cite{redmon2016yolo}.
Namun, deployment model-model ini pada perangkat edge dengan sumber
daya komputasi terbatas tetap menjadi tantangan signifikan, terutama
karena tingginya kebutuhan memori dan throughput komputasi dari
model berperforma tinggi \cite{jacob2018quantization}.

Teknik kompresi model, terutama kuantisasi bobot dan aktivasi,
telah terbukti efektif dalam mengurangi ukuran model dan meningkatkan
kecepatan inferensi tanpa kehilangan akurasi yang signifikan \cite{han2015deep}.
Dua pendekatan utama yang banyak diteliti adalah
\textit{Post-Training Quantization} (PTQ) dan
\textit{Quantization-Aware Training} (QAT).
PTQ mengkuantisasi model setelah pelatihan selesai menggunakan
dataset kalibrasi kecil, sehingga tidak memerlukan proses pelatihan
ulang \cite{banner2019post}.
Sebaliknya, QAT mengintegrasikan efek kuantisasi selama proses
pelatihan sehingga model dapat beradaptasi terhadap error kuantisasi,
menghasilkan akurasi yang lebih tinggi terutama pada presisi rendah \cite{jacob2018quantization}.

Meskipun perbandingan PTQ dan QAT telah banyak diteliti, terdapat
kesenjangan literatur yang signifikan pada kombinasi hardware tertentu.
Studi yang ada umumnya berfokus pada platform x86 CPU atau GPU NVIDIA \cite{wangqyolo2023},
smartphone berbasis ARM \cite{idama2024qatfp}, atau Raspberry Pi \cite{boddu2025yolov4}.
Sejauh penelusuran yang dilakukan, belum terdapat studi yang
secara sistematis membandingkan PTQ dan QAT pada arsitektur YOLO
modern di platform Apple Silicon M4 dengan Neural Engine,
GPU NVIDIA RTX 4060, dan perangkat MCU-class ESP32-S3 secara
bersamaan dalam satu kerangka eksperimen yang unified.

\section{Rumusan Masalah}

Berdasarkan latar belakang di atas, rumusan masalah penelitian ini adalah:
\begin{enumerate}
  \item Bagaimana trade-off latensi, throughput, akurasi, dan konsumsi daya
    antara PTQ dan QAT pada arsitektur YOLOv8n dan YOLOv26n di platform Apple Silicon M4?
  \item Bagaimana perbandingan karakteristik inferensi PTQ dan QAT antara
    platform Apple Silicon M4, GPU NVIDIA RTX 4060, dan MCU-class ESP32-S3?
  \item Format kuantisasi mana (FP32, FP16, PTQ INT8, QAT INT8) yang memberikan
    keseimbangan optimal antara akurasi dan efisiensi inferensi pada masing-masing platform?
\end{enumerate}

\section{Tujuan Penelitian}

\begin{enumerate}
  \item Mengukur dan membandingkan latensi, throughput, akurasi, dan konsumsi daya
    PTQ dan QAT pada YOLOv8n dan YOLOv26n di tiga platform hardware heterogen.
  \item Mengidentifikasi karakteristik hardware-specific yang mempengaruhi
    efektivitas masing-masing teknik kuantisasi.
  \item Memberikan panduan berbasis data empiris untuk pemilihan teknik
    optimasi inferensi pada sistem edge AI.
\end{enumerate}

\section{Batasan Penelitian}

\begin{enumerate}
  \item Model dibatasi pada YOLOv8n dan YOLOv26n.
  \item Platform hardware: Mac Mini M4, NVIDIA RTX 4060, dan ESP32-S3 Sense.
  \item Framework inferensi: ONNX Runtime, OpenVINO, TensorFlow Lite, dan CoreML.
  \item Dataset evaluasi: subset COCO val2017 dan dataset RescueVision.
  \item Pruning dan knowledge distillation berada di luar cakupan penelitian ini.
\end{enumerate}

\section{Manfaat Penelitian}

\begin{enumerate}
  \item Referensi empiris bagi pengembang sistem edge AI dalam memilih
    teknik optimasi inferensi sesuai karakteristik hardware.
  \item Benchmark baseline pada platform Apple Silicon M4 dan ESP32-S3
    yang belum tersedia dalam literatur yang ada.
  \item Pipeline benchmarking yang reproducible dan open-source
    sebagai kontribusi komunitas.
\end{enumerate}

\section{Sistematika Penulisan}

Skripsi ini disusun dalam lima bab. Bab I memaparkan latar belakang,
rumusan masalah, tujuan, batasan, dan manfaat penelitian.
Bab II mengulas landasan teori dan kajian pustaka terkait kuantisasi
model, arsitektur YOLO, dan platform edge computing.
Bab III menjelaskan metodologi penelitian meliputi desain eksperimen,
pipeline benchmark, dan prosedur pengukuran.
Bab IV menyajikan hasil eksperimen dan analisis komparatif.
Bab V memuat kesimpulan dan saran untuk penelitian selanjutnya.

% ════════════════════════════════════════════════════════════════════════════
\chapter{TINJAUAN PUSTAKA}
% ════════════════════════════════════════════════════════════════════════════

\section{Deteksi Objek Real-Time}

% TODO: ~1 halaman
% Poin: two-stage vs one-stage, YOLO sebagai dominant one-stage detector,
% evolusi dari v1 ke v11/v26, trade-off speed vs accuracy

\section{Kuantisasi Model Neural Network}

\subsection{Post-Training Quantization (PTQ)}

% TODO: ~1 halaman
% Poin: definisi, calibration dataset, MinMax vs percentile vs MSE,
% symmetric vs asymmetric, kelebihan (no retraining) vs kekurangan (accuracy drop)

\subsection{Quantization-Aware Training (QAT)}

% TODO: ~1 halaman
% Poin: fake quantization nodes, STE (Straight-Through Estimator),
% kelebihan (lebih akurat) vs kekurangan (butuh retraining, GPU cost)

\subsection{Perbandingan PTQ dan QAT}

% TODO: tabel perbandingan, cite Jacob 2018 dan Zhang 2025

\section{Arsitektur YOLO}

\subsection{YOLOv8}
% TODO: anchor-free, CSP backbone, decoupled head, parameter count

\subsection{YOLOv26n}
% TODO: arsitektur terbaru, perbandingan parameter dengan v8n

\section{Platform Edge Computing}

\subsection{Apple Silicon M4}
% TODO: Neural Engine, unified memory, CoreML, ONNX Runtime Metal EP

\subsection{NVIDIA RTX 4060}
% TODO: CUDA, TensorRT, ONNX Runtime CUDA EP

\subsection{ESP32-S3}
% TODO: Xtensa LX7, vector instructions, PSRAM, TFLite Micro

\section{Framework Inferensi}
% TODO: ONNX Runtime, OpenVINO, TFLite, CoreML — satu paragraf each

\section{Penelitian Terkait}

\begin{table}[H]
  \centering
  \caption{Perbandingan Penelitian Terkait}
  \label{tab:related_work}
  \begin{tabularx}{\textwidth}{lXXXX}
    \toprule
    \textbf{Penelitian} & \textbf{Model} & \textbf{Hardware} &
    \textbf{Teknik} & \textbf{Gap} \\
    \midrule
    \cite{idama2024qatfp} & Custom YOLO & Smartphone (Pixel 4) &
      QAT + Pruning & Tidak ada PTQ; model custom \\
    \cite{wangqyolo2023} & YOLOv5/v7 & RTX 4090, x86 CPU &
      PTQ & Tidak ada QAT; tidak ada Apple Silicon \\
    \cite{seo2025lltq} & YOLOv7/9/10 & PASCAL VOC only &
      QAT & Tidak ada deployment nyata \\
    \cite{boddu2025yolov4} & YOLOv4-Tiny & Raspberry Pi 5 &
      PTQ & Model lama; satu hardware; tidak ada QAT \\
    \textbf{Penelitian ini} & YOLOv8n, YOLOv26n &
      M4 + RTX 4060 + ESP32-S3 &
      PTQ vs QAT & Apple Silicon M4 + MCU-class; head-to-head \\
    \bottomrule
  \end{tabularx}
\end{table}

% ════════════════════════════════════════════════════════════════════════════
\chapter{METODOLOGI PENELITIAN}
% ════════════════════════════════════════════════════════════════════════════

\section{Desain Penelitian}

Penelitian ini menggunakan pendekatan eksperimental kuantitatif dengan
desain benchmarking sistematis. Variabel bebas adalah teknik kuantisasi
(FP32, FP16, PTQ INT8, QAT INT8) dan platform hardware (M4, RTX 4060, ESP32-S3).
Variabel terikat adalah latensi inferensi, throughput, akurasi deteksi,
ukuran model, dan konsumsi daya.

\section{Dataset}

\subsection{Dataset Kalibrasi PTQ}

Dataset kalibrasi PTQ menggunakan 500 gambar yang diambil secara acak
dari COCO train2017. Ukuran ini dipilih berdasarkan rekomendasi dari
literatur yang menunjukkan bahwa 100--500 gambar cukup untuk kalibrasi
PTQ yang stabil \cite{banner2019post}.

\subsection{Dataset Evaluasi Akurasi}

Evaluasi akurasi menggunakan dua dataset:
\begin{enumerate}
  \item COCO val2017 (5.000 gambar) untuk perbandingan dengan literatur yang ada.
  \item Dataset RescueVision (domain-specific rescue imagery) untuk evaluasi
    pada use case yang relevan.
\end{enumerate}

\section{Model}

\subsection{YOLOv8n}

YOLOv8n digunakan sebagai baseline yang well-established dengan referensi
literatur yang luas, memungkinkan perbandingan langsung dengan studi sebelumnya.

\subsection{YOLOv26n}

YOLOv26n merepresentasikan generasi terbaru arsitektur YOLO.
Baseline model telah tersedia dari pipeline RescueVision dengan
mAP@0.5 = 0.528 dan ukuran ONNX 11.70MB.

\section{Pipeline Eksperimen}

Pipeline eksperimen terdiri dari empat tahap utama:
\begin{enumerate}
  \item \textbf{Export}: Model PyTorch diekspor ke format ONNX sebagai representasi intermediate.
  \item \textbf{Optimasi}: Dari ONNX, model dioptimasi ke format target
    (FP16, PTQ INT8, QAT INT8) menggunakan framework yang sesuai per platform.
  \item \textbf{Benchmark}: Inferensi dijalankan dengan 10 iterasi warmup
    diikuti 100 iterasi pengukuran per konfigurasi.
  \item \textbf{Logging}: Semua hasil dicatat ke MLflow untuk reproducibility.
\end{enumerate}

\section{Prosedur Benchmark}

\subsection{Pengukuran Latensi}

Latensi diukur menggunakan \texttt{time.perf\_counter()} dengan input dummy
berukuran $640 \times 640 \times 3$. Warmup sebanyak 10 iterasi dilakukan
sebelum pengukuran untuk menghilangkan efek cold start.
Statistik yang dicatat: mean, standar deviasi, p95, p99.

\subsection{Pengukuran Akurasi}

Akurasi diukur menggunakan mAP@0.5 dan mAP@0.5:0.95 pada dataset evaluasi
dengan confidence threshold 0.25 dan IoU threshold 0.45.

\subsection{Pengukuran Konsumsi Daya}

Konsumsi daya diukur menggunakan:
\begin{itemize}
  \item Mac Mini M4: \texttt{powermetrics} (macOS)
  \item RTX 4060: \texttt{nvidia-smi} power query
  \item ESP32-S3: USB inline power meter
\end{itemize}

\section{Experiment Tracking}

Seluruh eksperimen ditracking menggunakan MLflow dengan parameter logging
meliputi: nama model, format kuantisasi, platform hardware, framework inferensi,
dan seluruh metrik hasil pengukuran. Dataset dan model checkpoint
di-version menggunakan DVC.

% ════════════════════════════════════════════════════════════════════════════
\chapter{HASIL DAN PEMBAHASAN}
% ════════════════════════════════════════════════════════════════════════════

\section{Hasil Benchmark Mac Mini M4}

% TODO: isi setelah eksperimen selesai

\begin{table}[H]
  \centering
  \caption{Hasil Benchmark YOLOv8n pada Mac Mini M4}
  \label{tab:m4_yolov8n}
  \begin{tabular}{lrrrr}
    \toprule
    \textbf{Format} & \textbf{Latency mean (ms)} &
    \textbf{p99 (ms)} & \textbf{FPS} & \textbf{mAP@0.5} \\
    \midrule
    FP32     & -- & -- & -- & -- \\
    FP16     & -- & -- & -- & -- \\
    PTQ INT8 & -- & -- & -- & -- \\
    QAT INT8 & -- & -- & -- & -- \\
    \bottomrule
  \end{tabular}
\end{table}

\begin{table}[H]
  \centering
  \caption{Hasil Benchmark YOLOv26n pada Mac Mini M4}
  \label{tab:m4_yolov26n}
  \begin{tabular}{lrrrr}
    \toprule
    \textbf{Format} & \textbf{Latency mean (ms)} &
    \textbf{p99 (ms)} & \textbf{FPS} & \textbf{mAP@0.5} \\
    \midrule
    FP32     & -- & -- & -- & -- \\
    FP16     & -- & -- & -- & -- \\
    PTQ INT8 & -- & -- & -- & -- \\
    QAT INT8 & -- & -- & -- & -- \\
    \bottomrule
  \end{tabular}
\end{table}

\section{Hasil Benchmark RTX 4060}
% TODO: isi setelah eksperimen selesai

\section{Hasil Benchmark ESP32-S3}
% TODO: isi setelah eksperimen selesai

\section{Analisis Komparatif Cross-Platform}
% TODO: analisis setelah semua data terkumpul

\section{Pembahasan}
% TODO: explain WHY angkanya seperti itu — ini bagian terpenting skripsi

% ════════════════════════════════════════════════════════════════════════════
\chapter{KESIMPULAN DAN SARAN}
% ════════════════════════════════════════════════════════════════════════════

\section{Kesimpulan}
% TODO: isi setelah penelitian selesai

\section{Saran}
% TODO: isi setelah penelitian selesai

% ─── Bibliography ────────────────────────────────────────────────────────────
\bibliographystyle{IEEEtran}
\bibliography{references}

\end{document}
